\chapter{Wiadomości wstępne.}
Regularnie będziemy aktualizowali ten ustęp o pojęcia z zewnątrz, które są potrzebne, a o których naturze miałem braki.

\section{Ciągi, szeregi, granice- nie tylko liczb.}

Wprowadzimy najpierw używaną terminologię, tzn. niech $A_{n}$ będzie \textbf{ciągiem zbiorów} z ustalonej przestrzeni X.
\begin{definition}[Ciąg rosnący(wstępujący) zbiorów.]

		Jeżeli dla każdego $n$ $A_{n}\subseteq A_{n+1} $ to mówimy, że ciąg $A_{n}$ jest \textbf{rosnący(wstępujący)}. 
\end{definition} 
\begin{definition}[Ciąg malejący(zstępujący) zbiorów.]
	
		Jeżeli dla każdego $n$ $A_{n} \supseteq A_{n+1}$ to mówimy, że ciąg $A_{n}$ jest \textbf{malejący(zstępujący)}. 	
\end{definition} 

\subsection{Czym są $\sum_{n=1}^{\infty}$ oraz $\prod_{n=1}^{\infty}$?}

W analizie matematycznej symbol $\sum_{n=1}^{\infty} \varphi\left( n \right) $ nie oznacza dodawania nieskończenie wielu liczb(bo nikt docześnie nie ma tyle czasu). Tak samo symbol $\prod_{n=1}^{\infty} \psi\left( n \right)$ nie oznacza mnożenia nieskończenie wielu liczb

Są to \textbf{skróty myślowe} oznaczające granicę \textbf{ciągu sum częściowych} oraz granicę \textbf{ciągu produktów częściowych} .

\textbf{Definicja:}
\begin{align*}
	\sum_{n=1}^{\infty} \varphi\left( n \right)  &\overset{def.}{=} \lim_{N \to \infty} \left( \sum_{n=1}^{N} \varphi\left( n \right)  \right) \\
	\prod_{n=1}^{\infty} \psi\left( n \right) &\overset{\Def}{=} \lim_{N \to \infty} \left( \prod_{n=1}^{N} \psi\left( n \right) \right)  
.\end{align*}

Nasuwa się tu analogia do tematyki \textbf{ciągu zbiorów}.

Mianowicie tak jak nieskończona suma(produkt) liczb, jest granicą ciągu sum częściowych, tak nieskończona suma(produkt) zbiorów, jest granicą ciągu ,,sum(produktów) częściowych zbiorów''.

\begin{enumerate}[label=\arabic*.]
	\item \textbf{Analogia: Szeregi Liczbowe vs. Szeregi Zbiorów} 
		\begin{itemize}
				\item \textbf{Świat Liczb(Szeregi):} Mamy ciąg liczb $\varphi_{n}$. Definiujemy \textbf{ciąg sum częściowych} $s_{N}$ :
			\[
			s_{N} = \sum_{n = 1}^{N} \varphi_{n} = \varphi_1 + \varphi_2 + \ldots + \varphi_{N}
			.\] 
			Wtedy suma nieskończona to granica sum częściowych:
			\[
			\sum_{n = 1}^{\infty} \varphi_{n} \overset{\Def}{=} \lim_{N\to \infty}\sum_{n = 1}^{N} \varphi_{n} 
			.\] 
				\item \textbf{Świat Zbiorów(Sumy mnogościowe):} Mamy ciąg zbiorów $A_{n}$. Definiujemy \textbf{ciąg sum częściowych} $S_{N}:$
			\[
			S_{N} = \bigcup_{n=1} ^{N} = A_1\cup A_2\cup \ldots\cup A_{N}
			.\] 
			Zauważmy kluczową rzecz: ten ciąg zbiorów jest \textbf{rosnący(wstępujący)}:
			\[
			S_1\subseteq S_2\subseteq S_3\subseteq \ldots
			.\] 
			Zatem analogicznie do liczb:
			\[
			\bigcup_{n=1} ^{\infty}A_{n} \overset{\Def}{=} \lim_{N \to \infty} S_{N} = \lim_{N \to \infty} \bigcup_{n=1} ^{N}A_{n}
			.\] 
		\end{itemize}
	\item \textbf{Analogia: Produkty Liczbowe vs. Produkty Zbiorów} 
		\begin{itemize}
			\item \textbf{Świat Liczb(Produkty):} Mamy ciąg liczb $\psi_{n}$. Definiujemy \textbf{ciąg produktów częściowych} $p_{N}$ :
			\[
			p_{N} = \prod_{n = 1}^{N} \psi_{n} = \psi_1 \cdot  \psi_2 \cdot  \ldots \cdot  \psi_{N}
			.\] 
			Wtedy produkt nieskończony to granica produktów częściowych:
			\[
			\prod_{n = 1}^{\infty} \psi_{n} \overset{\Def}{=} \lim_{N\to \infty}\prod_{n = 1}^{N} \psi_{n} 
			.\] 
				\item \textbf{Świat Zbiorów(Produkty(Iloczyny) mnogości owe)} Mamy ciąg zbiorów $A_{n}$. Definiujemy \textbf{ciąg produktów częściowych} $P_{N}:$
			\[
			P_{N} = \bigcap_{n=1} ^{N} = A_1\cap A_2\cap \ldots\cap A_{N}
			.\] 
			Zauważmy kluczową rzecz: ten ciąg zbiorów jest \textbf{malejący(zstępujący)}:
			\[
			P_1\supseteq P_2\supseteq P_3\supseteq \ldots
			.\] 
			Zatem analogicznie do liczb:
			\[
			\bigcap_{n=1} ^{\infty}A_{n} \overset{\Def}{=} \lim_{N \to \infty} P_{N} = \lim_{N \to \infty} \bigcap_{n=1} ^{N}A_{n}
			.\] 
 
		\end{itemize}

\end{enumerate}
\begin{definition}[Granica górna i granica dolna ciągu zbiorów $A_{n}$]
	Dla \textbf{ciągu zbiorów}  $\left( A_{n} \right)_{n} $ zbiory
	\[
		\limsup_{n \to \infty} A_{n} = \bigcap_{n=1} ^{\infty}\bigcup_{k=n} ^{\infty}A_{k}, \quad \liminf_{n \to \infty}A_{n} = \bigcup_{n=1}^{\infty}\bigcap_{k=n} ^{\infty}A_{k} 
	.\] 
	nazywamy odpowiednio \textbf{granicą górną} oraz \textbf{granicą dolną} ciągu $\left( A_{n} \right)_{n}$  
\end{definition} 
\begin{intuition}
	Jest to język do opisywania ,,zachowania granicznego'' punktów. Trzeba abstrahować od suchych znaczków i myśleć co one oznaczają.

	Wyobraźmy sobie $x$ oraz ciąg zbiorów $A_{n}$. Pytamy ,,co się dzieje z x w tym ciągu?''

	\begin{itemize}
		\item $\textbf{Granica Dolna(liminf): } \bigcup_{n=1}^{\infty}\bigcap_{k=n} ^{\infty}A_{k}$ 

			$\liminf_{n \to \infty} A_{n} = \lim_{n \to \infty} \left( \inf_{k\ge n} A_{k}  \right)= \sup \{ \inf\{A_{k}: k>n\}: n>0 \} $
			\begin{itemize}[label=$\ast$]
				\item $\bigcap_{k\ge n} A_{k}$ to zbiór punktów, które są we wszystkich $A_{n}$  od indeksu n wzwyż.
				\item $\bigcup_{n} \left( \ldots \right) $ oznacza, że x musi być w którymś z tych ogonów.
				\item \textbf{Intuicja:} $x\in \liminf_{n \to \infty} A_{n}$ oznacza, że x należy do $A_{n}$ dla \textbf{wszystkich $n$ poza conajwyżej skończoną liczbą}. Od pewnego miejsca $x$ ,,stabilizuje się'' i już zawsze jest w $A_{n}$.  
			\end{itemize}
		\item $\textbf{Granica Górna): } \bigcap_{n=1}^{\infty}\bigcup_{k=n} ^{\infty}A_{k}$ 

			$\limsup_{n \to \infty} A_{n} = \lim_{n \to \infty} (\sup_{k\ge n}A_{k}) = \inf \{ \sup \{A_{k}: k>n\}: n>0 \} $
			\begin{itemize}[label=$\ast$]
				\item $\bigcup_{k\ge n} $ to zbiór punktów, które są w jakimkolwiek $A_{k}$ od indeksu n wzwyż.
				\item $\bigcap_{n} \left( \ldots \right) $ oznacza, że x musi być w każdej takiej ,,ogonowej'' sumie. Nieważne jak daleko pójdziesz(jak duże n), x zawsze się jeszcze gdzieś pojawi w $A_{k}$ dla $k\ge n$.
				\item \textbf{Intuicja:}$x\in \limsup A_{n}$ oznacza, że x należy do $A_{n}$ dla nieskończnie wielu  $n$. x,,powraca'' bez końca.
			\end{itemize}
	\end{itemize}
	Zawsze $\liminf_{n \to \infty} A_{n} \subseteq \limsup_{n \to \infty} A_{n}$. Jeśli x jest prawie we wszystkich, to jest w niekończnie wielu). Ciąg $A_{n}$ jest zbieżny do $A$ jesli $\limsup_{n \to \infty}A_{n} = \liminf_{n \to \infty} A_{n}$.
\end{intuition} 

\begin{theorem}[Rozkład zbiorów otwartych.]
	Każdy niepusty zbiór otwarty $U\in \R$ jest postaci
	\[
	U = \bigcup_{n=1} ^{\infty}\left( a_{i}, b_{i} \right) 
	,\] 

	dla pewnych liczb wymiernych $a_{n}, b_{n}$
\end{theorem} 
\begin{proof}
	Każdego $x\in U$, możemy otoczyć liczbami wymiernymi- wynika to z gęstości liczb wymiernych.
\end{proof} 
\begin{remark}
	Mamy nieprzeliczalnie wiele punktów, ale wciskamy je w przeliczalnie wiele ,,klatek''. W pewien sposób ,,skompresowaliśmy'' nieprzeliczalny zbiór w przeliczalną sumę przedziałów.
\end{remark} 

Ważna własność odcinków domkniętych- w topologii nazywana \textbf{zwartością}. Nazywa też się to jakoś Twierdzeniem Heinego-Borela- z kazdego pokrycia zbioru domknietego mozna wybrac podpkrycie skonczone? Anyway tutaj w formie przedziałów:
\begin{theorem}[Pokrycie skończone przedziału domkniętego.]
	Jeżeli $[a,b] \subseteq \bigcup_{n=1} ^{\infty}(a_{n}, b_{n}) $ to istnieje $n\in N$, takie że $[a,b] \subseteq \bigcup_{i=1} ^{n}(a_{i},b_{i})$ 
\end{theorem} 
\begin{corollary}[Uogólnienie na zbiory zwarte(czyli ograniczone i domknięte)]
	Niech $F$ będzie domkniętym i ograniczonym podzbiorem prostej. Jeżeli  $F\subseteq \bigcup_{n=1} ^{\infty}(a_{i},b_{i})$ to istnieje $n\in \N$, takie że $F\subseteq \bigcup_{i=1} ^{n}(a_{i},b_{i})$.
\end{corollary}
\begin{intuition}
	Każdy zbiór domknięty i ograniczony ,,siedzi'' wewnątrz jakiegoś dużego odcinka $[A,B].$ Zbiór $F$ ma ,,dziury'' (tam gdzie go nie ma). Te dziury tworzą zbiór otwarty $F^{c}$.

	Konkretnie, to
	\begin{enumerate}[label=\arabic*.]
		\item Przykrywamy $F$ naszymi odcinkami otwartymi(w których $F$ się zawiera).
		\item Resztę dużego odcinka $[A,B]$ przykrywamy zbiorem  $F^{c}$, który jest otwarty.
		\item Teraz mamy pokrycie całego $[A,B]$.
		\item Z Tw. Heinego-Borela możemy wybrać dla niego podpokrycie skończone.
		\item  Z tego podpokrycia skończonego wyrzucamy $F^{c}$.
		\item Zostaje nam skończone pokrycie $F$.
	\end{enumerate}
\end{intuition} 
